\chapter{Varicocele}

\section*{Definition}
A varicocele is an abnormal dilation of the pampiniform venous plexus within the spermatic cord, caused by retrograde blood flow through incompetent valves in the internal spermatic vein. It is the most common correctable cause of male infertility, found in 15--20\% of all men and 35--40\% of men with primary infertility.

\section*{What Happens in Your Body}
Varicoceles are predominantly left-sided (85--90\%) due to the longer left testicular vein inserting into the left renal vein at a right angle. Venous reflux causes blood pooling, elevating intrascrotal temperature by 0.6--2\textdegree{}C above normal. This hyperthermia impairs spermatogenesis, induces oxidative stress in testicular tissue, and damages sperm DNA integrity.

\section*{Causes}
\begin{itemize}
  \item \textbf{Primary (unknown cause):} Most common; caused by incompetent or congenitally absent venous valves in the internal spermatic vein
  \item \textbf{Secondary:} Compression of the left renal vein (nutcracker syndrome between the aorta and superior mesenteric artery), retroperitoneal tumor, hydronephrosis, situs inversus, renal cell carcinoma with tumor thrombus, liver scarring (portosystemic collaterals)
  \item \textbf{Risk factors:} Tall, lean body habitat (longer venous column increases hydrostatic pressure), family history, prolonged standing
\end{itemize}

\section*{When to See a Doctor}
See your doctor if:
\begin{itemize}
  \item A scrotal mass or "bag of worms" sensation is noted, particularly when standing
  \item Dull, aching scrotal pain that worsens with standing or exertion and improves with lying down
  \item Infertility after 12 months of unprotected intercourse with abnormal semen analysis
  \item Acute onset or right-sided varicocele (requires urgent evaluation for retroperitoneal pathology)
\end{itemize}

\section*{Related Symptoms}
\begin{itemize}
  \item Dull scrotal ache, testicular atrophy (especially left), infertility, scrotal heaviness, visible or palpable scrotal mass
\end{itemize}
\textbf{Important:} Most varicoceles are asymptomatic and found incidentally on physical exam; only 10--20\% of men with varicocele report pain.

\section*{Self-Care / Home Management}
\begin{itemize}
  \item Wear supportive underwear (briefs instead of boxers) to reduce venous pooling
  \item Avoid prolonged standing and heavy lifting
  \item Lie supine when pain is bothersome to improve venous drainage
  \item Apply cold compresses to the scrotum for pain relief
\end{itemize}

\section*{How Your Doctor Will Evaluate You}
\begin{itemize}
  \item Physical exam in standing position and with Valsalva maneuver; graded 1--3 (1: palpable only with Valsalva, 2: easily palpable, 3: visible through scrotal skin)
  \item Scrotal ultrasound with Doppler to confirm venous reflux, measure testicular size, and exclude secondary causes
  \item Semen analysis if infertility is a concern
  \item If right-sided or acute onset: CT abdomen/pelvis to exclude retroperitoneal mass
\end{itemize}

\section*{Treatment Options}
\begin{itemize}
  \item No treatment needed for asymptomatic varicoceles not impairing fertility
  \item Indications for intervention: ipsilateral testicular atrophy (catch-up growth), abnormal semen parameters with infertility, persistent pain, bilateral varicoceles
  \item Options: microsurgical varicocelectomy (gold standard, lowest recurrence), laparoscopic varicocele ligation, or percutaneous embolization
  \item Improvement in semen parameters in 50--70\% after repair; pregnancy rate ~40--50\%
  \item Post-procedure recurrence rate: 1--2\% for microsurgery, up to 15\% for non-microsurgical approaches
\end{itemize}

\section*{References}
\begin{thebibliography}{99}
\bibitem{aua-varicocele} Sharlip ID, et al. "AUA guideline for the management of varicocele." \emph{J Urol}. 2024;212(2):218-228.
\bibitem{uptodate-varicocele} Wein AJ, et al. Varicocele: clinical features and evaluation. In: UpToDate, Post TW (Ed), UpToDate, Waltham, MA; 2025.
\bibitem{nejm-varicocele} Schlegel PN, et al. "Varicocele and Male Infertility." \emph{N Engl J Med}. 2023;389(14):1305-1315.
\bibitem{bmj-varicocele} Clavijo RI, et al. "Varicocele in adults." \emph{BMJ}. 2024;386:e081543.
\bibitem{cochrane-varicocele} Kroese ACJ, et al. "Varicocele repair for male infertility." \emph{Cochrane Database Syst Rev}. 2023;(6):CD014838.
\end{thebibliography}
